\documentclass[11pt]{article}
\usepackage{geometry}                % See geometry.pdf to learn the layout options. There are lots.
\geometry{letterpaper}                   % ... or a4paper or a5paper or ... 
%\geometry{landscape}                % Activate for for rotated page geometry
\usepackage[parfill]{parskip}    % Activate to begin paragraphs with an empty line rather than an indent
\usepackage{graphicx}
\usepackage{amssymb}
\usepackage[normalem]{ulem}
\usepackage{epstopdf}
\usepackage{tcolorbox}
\usepackage{bm}
\usepackage{float}
\restylefloat{figure}
\usepackage{titlesec}
\usepackage{verbatimbox}
\usepackage{multicol}
\usepackage{soul}
\usepackage{tikz}
\usepackage{tikzsymbols}
\usepackage{amsmath}
\usetikzlibrary{arrows}
\usepackage{caption}
\usepackage{longtable}
\DeclareGraphicsRule{.tif}{png}{.png}{`convert #1 `dirname #1`/`basename #1 .tif`.png}
\pagestyle{plain}
%\setcounter{page}{1}
\renewcommand{\thepage}{R -- \arabic{page}} 
\setcounter{page}{1} 

\setlength{\textwidth}{6.9in}
\setlength{\oddsidemargin}{-.35in}
\setlength{\evensidemargin}{-.35in}
\setlength{\textheight}{9.2in}
%\setlength{\topmargin}{-0.5in}
\setlength{\topmargin}{-.9in}
%\setlength{\oddsidemargin}{-.25in}
%\setlength{\evensidemargin}{-.25in}
%\pagestyle{headandfoot}
%\firstpageheader{Stat 341}{Survey of Statistical Methods}{}

\newcommand{\ub}[1]{\ul{\bf #1}}

\newcommand{\vsp}{\vspace{1in}}


%\title{Brief Article}
%\author{The Author}
%\date{}                                           % Activate to display a given date or no date

\begin{document}
\large
\begin{tcolorbox}[width=\linewidth, sharp corners=all, colback=white!100!black]

%%% TOPIC 1

\begin{center} {\huge{\bf R Reference Sheet}} \end{center}

\end{tcolorbox}

%\begin{tcolorbox}[width=\linewidth, sharp corners=all, colback=white!95!black, title = {\bf Topic 1}, colbacktitle = white!80!black, coltitle = black]
%\centerline{{\Large\bf Data and Variables}}
%\end{tcolorbox}

Tips:\vspace{-.1in}
\begin{enumerate}
	\item When in the console, the up arrow retrieves your previous command.
	\item You can auto-complete by hitting the tab key.
	\item We always need to specify the path to a file unless it is in the working directory. We can also specify the path from the working directory to the file.
	\item Typing \texttt{Control+Enter} on Windows or \texttt{Command+Return} on Mac will run the line of code your cursor is on.
	\item Typing \texttt{Option+minus} will enter the arrow operator \texttt{ <- } with spaces around it.
	\item Typing \texttt{Control+Shift+M} on Windows or \texttt{Command+Shift+M} gives the pipe (\texttt{|>} or \texttt{\%>\%}).
	\item In {\bf R}, the comment symbol is \texttt{\#}.
	\item Type a question mark in front of a function or data set to get documentation about it. Example: \texttt{?rep}.
\end{enumerate}

%\vspace{.1in}
\begin{tcolorbox}[width=2.75in, sharp corners=all, colback=white!95!black] %title = {\bf Day 1}, colbacktitle = white!80!black, coltitle = black]
{\bf Important Base R Functions}
\end{tcolorbox}

\vspace{-.1in}

\begin{longtable}{clcc}
	\underline{\bf Function} & \underline{\bf Explanation}\\
	\texttt{c()} & Creates a vector.\\
	\texttt{data.frame()} & Creates a data frame.\\
	\texttt{install.packages("")} & Installs the indicated package. The name must be in quotes.\\
	& Only need to use once per package.\\
	\texttt{library()} & Loads a library/package to give access to its functions and data sets.\\
	\texttt{read.csv("")} & Reads in a comma separated value file. Put file name in quotes.\\
	\texttt{read.table("")} & Reads in a text file. Put file name in quotes.\\
	\texttt{factor()} & Turns a vector into a factor (categorical variable).\\
	\texttt{length()} &  Counts the number of elements in a vector.\\
	\texttt{rep()} & Repeat a value, string, or vector.\\
	\texttt{mean()} & Computes the average of a vector.\\
	\texttt{sd()} & Computes the standard deviation of a vector. \\
	\texttt{median()} & Computes the median of a vector.\\
	\texttt{quantile()} & Computes a quantile. Put 0.25 as second entry for $Q_1$, 0.75 for $Q_3$.\\
	\texttt{min()} & Finds the minimum of a vector.\\
	\texttt{max()} & Finds the maximum of a vector.\\
	\texttt{sum()} & Adds up all values in a vector.\\
	\texttt{table()} & Counts how many of each category or value for a vector.\\
	\texttt{summary()} & Computes the 5-number summary (and the mean) of a vector.\\
	\texttt{round()} & Rounds a number or vector to the desired number of decimal places.\\
	\texttt{data()} & Loads a data set into memory.\\
	\texttt{class()} & Find the class an object belongs to, like vector, factor, data frame, etc.\\
	\texttt{exp()} & Compute the number $e$ raised to a power.\\
	\texttt{log()} & Compute the log of a number (natural log by default).\\
	\texttt{unique()} & Returns all of the different (unique) values in a vector. \\
	\texttt{cbind()} & Binds columns.\\
	\texttt{rbind()} & Binds rows.\\
	\texttt{nrow()} & Gives number of rows in data frame.\\
	\texttt{ncol()} & Gives number of columns in data frame.\\
	\texttt{dim()} & Gives number of rows and columns in data frame.\\
	\texttt{rm()} & Removes a variable from the environment.\\
	\texttt{as.factor()} & Changes the type of a vector to a factor.\\
	\texttt{as.character()} & Changes the type of a vector to a character.\\
	\texttt{as.numberic()} & Changes the type of a vector to numeric.\\
	\texttt{na.omit()} & Removes all NA values from a vector or all NA rows from a data frame.\\
	\texttt{function()\{ \}} & Creates a function. The function body should be between \texttt{\{} and \texttt{\}}.
\end{longtable}


%\vspace{.1in}
\begin{tcolorbox}[width=3.2in, sharp corners=all, colback=white!95!black] %title = {\bf Day 1}, colbacktitle = white!80!black, coltitle = black]
	{\bf Assigning and Subsetting Vectors}
\end{tcolorbox}

\begin{longtable}{lcll}
	\underline{\bf Command} & \underline{\bf Result} & & \underline{\bf Explanation}\\
	\texttt{x <- c(2, 4, 6, 8)} & \texttt{2 4 6 8}  && Join elements into a vector.\\
	\texttt{x = c(2, 4, 6, 8)} & \texttt{2 4 6 8} &&  Alternative way to join elements into a vector.\\
	\texttt{2:6} & \texttt{2 3 4 5 6} & & An integer sequence.\\
	\texttt{seq(2, 3, by = 0.5)} & \texttt{2.0 2.5 3.0} && A more complicated sequence.\\
	\texttt{rep(1:2, times = 3)} & \texttt{1 2 1 2 1 2} && Repeat a vector.\\
	\texttt{rep(1:2, each = 3)} & \texttt{1 1 1 2 2 2} && Repeat elements of a vector.\\
	\texttt{x[4]} & \texttt{8} && Take the fourth element of \texttt{x}.\\
	\texttt{x[-4]} & \texttt{2 4 6} && Take all but the fourth element of \texttt{x}.\\
	\texttt{x[2:4]} & \texttt{4 6 8} && Take elements two through four of \texttt{x}.\\
	\texttt{x[c(1, 3)]} & \texttt{2 6} && Take elements one and three of \texttt{x}.\\
	\texttt{x[-c(2, 3)]} & \texttt{2 8} && Take all but the second and third elements of \texttt{x}.\\
	\texttt{x[x == 6]} & \texttt{6} && Take elements that are equal to 10 in \texttt{x}.\\
	\texttt{x[x < 5]} & \texttt{2 4} && Take elements that are less than 5 in \texttt{x}.\\
	\texttt{x[x \%in\% c(1, 2, 3, 4)]} & \texttt{2 4} && Take elements in the set 1, 2, 3, 4 in \texttt{x}.\\
	\texttt{x["apple"]} &  && Take elements in \texttt{x} named ``apple''.
\end{longtable}


%\vspace{.1in}
\begin{tcolorbox}[width=4.4in, sharp corners=all, colback=white!95!black] %title = {\bf Day 1}, colbacktitle = white!80!black, coltitle = black]
	{\bf Assigning and Subsetting Data Frames or Tibbles}
\end{tcolorbox}


\begin{longtable}{llc}
	\underline{\bf Command} &  \underline{\bf Explanation}\\
	\texttt{df <- data.frame(}  & \\
	\phantom{.}\hspace{0.1in}\texttt{x = 1:3,} & Creates data frame with columns \texttt{x} and \texttt{y}.\phantom{move this over over over!}\\
	\phantom{.}\hspace{0.1in}\texttt{y = c("a", "b", "c")} & An equal sign can also be used instead of the arrow.\\
	\texttt{)} & \\
	\texttt{df\$x} & Takes the \texttt{x} column in the data frame.\\
	\texttt{df\$x[1:3]} & Takes the 1st through 3rd element the \texttt{x} column.\\
	\texttt{df[, 2]} & Takes everything in the 2nd column of the data frame.\\
	\texttt{df[3, ]} & Takes everything in the 3rd row of the data frame.\\
	\texttt{df[2:3, 1]} & Takes the 2nd and 3rd rows in the first column.\\	
	\texttt{df[[1]]} & Takes the first column.
\end{longtable}

\newpage


\begin{tcolorbox}[width=2.95in, sharp corners=all, colback=white!95!black] %title = {\bf Day 1}, colbacktitle = white!80!black, coltitle = black]
	{\bf Important Tidyverse Functions}
\end{tcolorbox}
Type \texttt{library(tidyverse)} after you've installed it with \texttt{install.packages("tidyverse")} to load the tidyverse and have access to these functions. 



\begin{longtable}{clcc}
	\underline{\bf Function} & \underline{\bf Explanation}\\
	\texttt{tibble()} & Creates a tibble, which is a more modern data frame.\\
	\texttt{mutate(data, variable)} & Creates a new variable in the given data frame.\\
	\texttt{filter(data, condition)} & Subsets data frame based on the condition. Selects rows.\\
	\texttt{select(data, variables)} & Selects variables based on their name. Selects columns.\\
	\texttt{summarize(data, functions)} & Computes the function value for a variable in the data frame.\\
	\texttt{arrange(data, variable)} & Changes the ordering of the rows. The \texttt{desc()} function is useful.\\
	\texttt{group\_by(data, variable)} & allows you to perform any operation ``by group".\\
	 & Useful before using \texttt{summarize()}.\\
	 \texttt{pivot\_longer(data, cols)} & Changes data format from not-tidy to tidy.\\
	 \texttt{read\_csv("")} & Reads in a csv file as a tibble. Put file name in quotes.\\
	 \texttt{read\_table("")} & Reads in a text file as a tibble. Put file name in quotes.\\
	\texttt{read\_excel("")} & Read in excel file. Put file name in quotes. Requires \texttt{readxl} package.
\end{longtable}


%\vspace{.1in}
\begin{tcolorbox}[width=2.8in, sharp corners=all, colback=white!95!black] %title = {\bf Day 1}, colbacktitle = white!80!black, coltitle = black]
	{\bf Important ggplot2 Functions}
\end{tcolorbox}
\texttt{ggplot2} is part of the tidyverse, so typing \texttt{library(tidyverse)} will also load \texttt{ggplot2}.

Every ggplot needs three components: 
\begin{enumerate}
	\item The data in tidy format (one observation/data value per row with variables in each column). 
	\item An aesthetic mapping using the \texttt{aes()} function, which usually goes in the \texttt{ggplot()} function. 
	\item A geometry using a \texttt{geom\_something()} functon.
\end{enumerate}


\begin{longtable}{clcc}
	\underline{\bf Function} & \underline{\bf Explanation}\\
	\texttt{ggplot(data, aes())} & The core plotting function. Needs a data frame/tibble entered and often \\
	& has the \texttt{aes()} function if the arguments of \texttt{aes()} apply to all.\\
	\texttt{aes()} & The aesthetic map. Typical arguments include \texttt{x}, \texttt{y}, \texttt{fill}, \texttt{color}, \texttt{shape}.\\
	\texttt{geom\_point()} & Creates a scatter plot. \\
	\texttt{geom\_bar()} & Creates a bar plot.  \\
	\texttt{geom\_boxplot()} & Creates a box plot. \\
	\texttt{geom\_histogram()} & Creates a histogram. \\
	\texttt{geom\_density()} & Creates a density plot. \\
	\texttt{geom\_abline()} & Adds a line to the plot with the \texttt{slope} and \texttt{intercept} arguments. \\
	\texttt{geom\_qq()} & Creates a quantile-quantile (QQ) plot. \\
	\texttt{geom\_qq\_line()} & Put a QQ line on a QQ plot. \\
	\texttt{labs()} & Changes the labels for  \texttt{x}, \texttt{y},  \texttt{title}, \texttt{fill}, \texttt{color}, \texttt{shape}, etc.\\
	\texttt{lims()} & Changes the limits on \texttt{x}, or \texttt{y} or both.\\
	\texttt{theme()} & Allows for the adjustments of many options on the graph, like text size.\\
	\texttt{theme\_bw()} & One of several \texttt{theme\_} functions that changes the plot style. Put this \\
	& function before the \texttt{theme()} function if both are used.
\end{longtable}

Each of the above functions, with the exception of \texttt{ggplot()} and \texttt{aes()}, is added on to the previous function with a plus sign (\texttt{+}). An example of some ggplot code is given below:

\begin{verbatim}
ggplot(mtcars, aes(x = wt, y = mpg, color = factor(cyl))) +
  geom_point(size = 2) +
  labs(x = "Weight", y = "MPG", color = "Cylinders", title = "ggplot!") +
  theme(axis.title = element_text(size = 14),
         plot.title = element_text(size = 18, face = "bold"),
         legend.position = "bottom")
\end{verbatim}


\vspace{.1in}
\begin{tcolorbox}[width=3.35in, sharp corners=all, colback=white!95!black] %title = {\bf Day 1}, colbacktitle = white!80!black, coltitle = black]
	{\bf Probability Functions}
\end{tcolorbox}

\begin{tabular}{clcc}
	\underline{\bf Function} & \underline{\bf Explanation}\\
	\texttt{dbinom()} & Finds $P(X=x)$ for a binomial. \\
	\texttt{pbinom()} & Finds $P(X\le x)$ for a binomial. \\
	\texttt{rbinom()} & Randomly generates values from a binomial distribution. \\
	\texttt{dpois()} & Finds $P(X=x)$ for a Poisson. \\
	\texttt{ppois()} & Finds $P(X\le x)$ for a Poisson. \\
	\texttt{rpois()} & Randomly generates values from a Poisson distribution. \\
	\texttt{pnorm()} & Finds $P(X\le x)=P(X<x)$ for a normal.\\
	\texttt{qnorm()} & Finds a percentile for a normal.\\
	\texttt{rnorm()} & Randomly generates values from a normal distribution.\\
\end{tabular}

{\bf R} also has functions that begin with \texttt{p},  \texttt{q}, and  \texttt{r} for  \texttt{t},  \texttt{chisq},  \texttt{f} that work the same way as the \texttt{pnorm()},  \texttt{qnorm()}, and  \texttt{rnorm()} functions.\\

\vspace{.1in}
\begin{tcolorbox}[width=3.35in, sharp corners=all, colback=white!95!black] %title = {\bf Day 1}, colbacktitle = white!80!black, coltitle = black]
	{\bf More Advanced Statistics Functions}
\end{tcolorbox}

\begin{tabular}{clcc}
	\underline{\bf Function} & \underline{\bf Explanation}\\
		\texttt{cor(x, y)} & Computes the correlation between two vectors.\\
		\texttt{t.test()} & Performs a t test for one or two means.\\
		\texttt{prop.test()} & Performs a test for one or two proportions.\\
		\texttt{chisq.test()} & Performs a goodness-of-fit test or test for independence.\\
		\texttt{cor.test()} & Performs a correlation test.\\
		\texttt{lm(y $\sim$ x, data = df)} & Creates a linear model.\\
		\texttt{glm(y $\sim$ x, data, family)} & Creates a generalized linear model.\\
		\texttt{aov(y $\sim$ x, data = df)} & Creates an ANOVA object.\\
		\texttt{anova(model)} & Outputs an ANOVA table for a given model.\\
		\texttt{predict()} & Obtain predictions, confidence intervals, and/or prediction intervals.\\
		\texttt{prcomp()} & Used for principal component analysis. \\
		\texttt{wilcox.test()} & Performs a Wilcoxon rank sum test. \\
		\texttt{kruskal.test()} & Performs a Kruskal-Wallis  test. \\
\end{tabular}





\end{document}   
